\section{Introduction}
\label{sec:intro}

Wildlife management is based on reliable estimates of wildlife abundance. White-tailed deer
are among the most intensively managed ungulates in North America, and their density informs
hunting quotas, culling decisions and vehicle-collision mitigation. The stakes extend beyond
harvest. Surveillance for chronic wasting disease is most efficient when it concentrates on
where infection risk is highest~\cite{jennelle2018applying}, and deer-to-human transmission of
SARS-CoV-2 was identified in animals sampled during routine surveillance of exactly that
kind~\cite{pickering2022divergent}. Each of these decisions depends on knowing how many
animals are present.

For larger species such as deer, one established way of estimating abundance is for
observers to drive transects at night with spotlights or handheld thermal scopes. Road
transects are practical but not neutral. Spotlight surveys along roads carry a known
bias~\cite{vercauteren2011managing}, and a road used as a transect line can violate the
assumption, central to distance sampling, that animals are distributed uniformly with respect
to the line, because deer may be drawn to or kept away from the road
itself~\cite{diefenbach2025accounting}. Thermal imaging extends the survey into darkness by
detecting emitted heat, but it brings constraints of its own: contrast falls as the ground
approaches body temperature, vegetation hides animals, and a distant deer occupies only a few
pixels~\cite{burke2019optimizing}. Any automated survey inherits all of these at once.

Automating the survey is an obvious application of modern object detection, and the
literature has largely treated it as one. From camera traps~\cite{norouzzadeh2018automatically}
to aerial imagery~\cite{kellenberger2018detecting}, systems are trained on annotated animal
imagery and judged by per-image accuracy, typically mean average precision (mAP), and the
survey problem is regarded as addressed once that number is high.

This paper argues, and then measures, that the two problems come apart. A detector is
rewarded for finding animals in \emph{frames}. A survey needs the number of distinct
\emph{animals}. Between those sits a chain of decisions (associating detections into tracks,
deciding which tracks are real, deciding which are duplicates of an animal already counted)
that no detection metric scores.

That distinction also fixes the yardstick we use throughout. A survey does not need a box
placed precisely on an animal; it needs the animal noticed once. We therefore score the
detection stage \emph{per animal}, over the whole of an animal's track, asking whether any
frame produced a box touching it, and report it alongside the conventional per-box metrics
rather than in place of them. Per-box mAP remains a valid measure of localization, but it
answers a question a survey does not ask, and read as a survey metric it can place the
difficulty in the wrong stage.

Which stage limits a counting system is also not something to assume. In sonar fish counting,
substituting a perfect detector recovers almost all of the counting error, so there detection
is the binding constraint~\cite{kay2022caltech}. A thermal deer transect differs in ways that
could move the constraint elsewhere: the animals are faint, but each is visible over many
consecutive frames as the vehicle passes, which gives association and confirmation far more
evidence to work with. Settling the question needs an instrument that attributes error to each
stage separately, applied to data annotated at the level of individual animals.

We study this on a corpus built for the purpose: 32 FLIR thermal road transects from four
sites in Illinois, 521,930 frames at $640\times512$ and 60\,fps, with every individual deer
annotated as a distinct track. Because one annotated track is one animal, the corpus supplies
detection, tracking and count ground truth simultaneously, the property on which the
decomposition below depends.

Our central instrument is a three-stage decomposition of counting error. For each
ground-truth animal we ask whether it is \reachedq (touched by at least one candidate track),
whether it is \primaryq (some candidate covers it better than it covers any other animal, so a
one-count-per-track decision \emph{could} count it), and whether it is finally \countedq. The
gaps between these three quantities attribute error to detection, to association and to
confirmation respectively.

With the loss localized, we test both remedies a practitioner would reach for. We enrich the candidate pool by lowering the detector's confidence threshold,
removing the tracker's track-initialization gate, relaxing non-maximum suppression (NMS), and
adding an appearance term to association. We then replace the rule-based confirmation stage
with learned alternatives of increasing capacity, under three evaluation protocols. In each
case the question is the same: does the count improve?

The paper contributes three things. The first is the corpus: fully annotated thermal road
transects with per-individual identity, which makes detection, tracking and counting
evaluable against one set of annotations, released with the evaluation harness. The second is
the decomposition, together with a held-out protocol that evaluates counting only on videos
the detector never trained on. The third is a controlled test of where counting error arises
and whether improving detection or confirmation reduces it. The remaining analyses (the
detector benchmark, the matching-criterion analysis, the sensitivity of the confirmation
rule, the limits of appearance matching at 27 pixels, and the transfer between sites) support
or qualify these three.
